\appchapter{Mã nguồn chương trình}
\label{appendix:source_code}

Toàn bộ mã nguồn của hệ thống được công khai tại:

\begin{center}
    \url{https://github.com/Huypham07/esg-washing}
\end{center}

Hai mô hình phân loại đã huấn luyện được publish trên HuggingFace và tự động tải về khi chạy pipeline:

\begin{itemize}
    \item \url{https://huggingface.co/huypham71/esg-topic-classifier}
    \item \url{https://huggingface.co/huypham71/esg-action-classifier}
\end{itemize}

\section*{Cấu trúc thư mục}

\begin{verbatim}
esg-washing/
├── config/
│   ├── pipeline.yml          # Tham số pipeline (EWRI, mô hình, ngưỡng)
│   └── train.yml             # Siêu tham số huấn luyện
├── data/
│   ├── extracted/            # Văn bản thô / ZIP source PDF
│   └── labels/               # Dữ liệu gán nhãn (topic, action)
├── notebooks/
│   ├── c4_experiments.py     # Thực nghiệm Chương 4 (RQ1/RQ2/RQ3)
│   └── train-model.ipynb     # Huấn luyện trên Colab (GPU)
├── outputs/
│   └── experiments/          # Kết quả liên kết bằng chứng, Grid Search
├── src/
│   ├── pipeline/
│   │   ├── run.py            # Điểm vào CLI — phân tích một báo cáo
│   │   ├── pipeline.py       # ESGWashingPipeline (điều phối toàn bộ)
│   │   ├── document_loader.py    # Tải PDF / TXT, OCR qua Docling
│   │   ├── evidence_detector.py  # Phát hiện loại bằng chứng (regex)
│   │   ├── evidence_linker.py    # Liên kết bằng chứng (TF-IDF + SBERT)
│   │   ├── nli_verifier.py       # Xác minh NLI (mDeBERTa)
│   │   ├── ewri.py               # Tính WRS và EWRI
│   │   ├── ewri_grid_search.py   # Grid Search tham số EWRI
│   │   └── evidence_experiments.py # Thực nghiệm 3 biến thể bằng chứng
│   └── training/
│       ├── train_model.py        # Huấn luyện neuro-symbolic / baseline
│       ├── neuro_symbolic.py     # Semantic loss (exactly-one, impl, neg)
│       ├── tune_hyperparams.py   # Tối ưu siêu tham số (Optuna)
│       ├── corpus/
│       │   └── build_corpus.py   # Tiền xử lý PDF → ngữ liệu câu
│       └── labeling/
│           ├── grounded_rules.py     # Bộ quy tắc ký hiệu (GRI/Bloom/Hyland)
│           ├── topic_llm_labeler.py  # Gán nhãn chủ đề bằng Gemini
│           └── action_llm_labeler.py # Gán nhãn hành động bằng Gemini
└── requirements.txt          # Dependencies Python
\end{verbatim}

\section*{Hướng dẫn chạy}

\subsection*{Cài đặt môi trường}

Yêu cầu Python 3.10 trở lên. Cài đặt dependencies:

\begin{verbatim}
pip install -r requirements.txt
\end{verbatim}

\subsection*{1. Phân tích một báo cáo thường niên (Demo)}

\begin{verbatim}
python -m src.pipeline.run \
  --input  path/to/annual_report.pdf \
  --output outputs/demo \
  --bank   "TEN_NGAN_HANG" \
  --year   2024
\end{verbatim}

Kết quả xuất ra thư mục \texttt{outputs/demo/}: tệp \texttt{report.html} (báo cáo HTML gồm điểm EWRI, phân rã thành phần, top-10 câu rủi ro cao nhất), \texttt{enriched.parquet} (dữ liệu mức câu đầy đủ), và \texttt{extracted.txt} (văn bản sau OCR).

\subsection*{2. Huấn luyện mô hình phân loại}

\begin{verbatim}
# Neuro-Symbolic (có semantic loss)
python -m src.training.train_model --task topic
python -m src.training.train_model --task action

# Baseline (không có semantic loss, dùng so sánh RQ1)
python -m src.training.train_model --task topic  --no-neuro-symbolic
python -m src.training.train_model --task action --no-neuro-symbolic
\end{verbatim}

Model được lưu vào \texttt{outputs/models/\{topic,action\}\_classifier/final/}.

\subsection*{3. Liên kết bằng chứng và Grid Search tham số EWRI}

\begin{verbatim}
# Bước 3a: chạy 3 biến thể liên kết bằng chứng trên toàn corpus
python -m src.pipeline.evidence_experiments \
  --input     data/corpus/actionability_sentences.parquet \
  --output-dir outputs/experiments/evidence

# Bước 3b: Grid Search tham số (P_action, lambda, C)
python -m src.pipeline.ewri_grid_search \
  --input       outputs/experiments/evidence/evidence_nli.parquet \
  --output      outputs/ewri_grid_search_results.csv \
  --figures-dir thesis/figures
\end{verbatim}

\subsection*{4. Thực nghiệm Chương 4 (RQ1/RQ2/RQ3)}

Mở \texttt{notebooks/c4\_experiments.py} trong Jupyter hoặc VS~Code (Jupyter mode). Notebook chia thành ba phần tương ứng ba câu hỏi nghiên cứu: RQ1 (đánh giá classifier), RQ2 (so sánh chiến lược bằng chứng), và RQ3 (phân tích EWRI trên 50 quan sát ngân hàng-năm). Biểu đồ được xuất tự động vào \texttt{thesis/figures/}.
